% !PRINTED 125
\weberpage{125}

Es ist nun ferner leicht einzusehen, dass $\gamma=g$ sein muss, denn da $g$ die \gesperrt{untere Grenze} aller Werthe von $X$ innerhalb $(R)$ ist, so kann $\gamma$ zunächst nicht kleiner als $g$ sein.

Es kann aber $\gamma$ auch nicht grösser als $g$ sein; denn alle Werthe, die $X$ auf der Strecke $MM'$ annimmt, sind $\ge\gamma$; ist aber $\gamma>g$, so kann man eine Grösse $g'$ so annehmen, dass $\gamma>g'>g$ ist, und wegen der in \S.\,31, VI bewiesenen Stetigkeit von $X$ lassen sich die Zahlen $\alpha,\beta$ so bestimmen, dass
\[
 \alpha<\eta<\beta
\]
und dass in dem ganzen Bereich $(P,\beta)-(P,\alpha)=(\alpha,\beta)$ ($QNN'Q'$ in der Fig. 2) $X$ grösser als $g'$ bleibt. Folglich ist die \gesperrt{untere Grenze} von $X$ in $(\alpha,\beta)$ grösser als $g$. Da nun die \gesperrt{untere Grenze} von $X$ in $(P,\alpha)$ grösser als $g$ ist, so ist sie auch in $(P,\beta)$, was aus $(P,\alpha)$ und $(\alpha,\beta)$ zusammengesetzt ist, grösser als $g$; $\beta$ gehört aber zu $\mathfrak B$, woraus ein Widerspruch mit der Definition von $\mathfrak B$ folgen würde.

Es bleibt also nur übrig, dass $\gamma=g$ ist, und der Satz 4. ist damit nachgewiesen, nämlich, dass es einen Werth $\xi$ im Inneren von $(R)$ giebt, für den
\[
 |f(\xi)|=g
\]
ist, dass also $|f(x)|$ einen Minimumwerth erreicht. Wir beweisen nun:

5. \gesperrt{Wenn $\alpha$ irgend ein Werth von $x$ ist, für den $f(\alpha)$ von Null verschieden ist, so lässt sich $h$ so annehmen, dass}
\begin{equation}
 |f(\alpha+h)|<|f(\alpha)|
\tag{10}
\end{equation}
\gesperrt{ausfällt.}

Daraus folgt dann, dass, wenn $f(\xi)$ von Null verschieden wäre, $g$ nicht das Minimum der Function $|f(x)|$ sein könnte; da dies aber bewiesen ist, so muss
\begin{equation}
 f(\xi)=0
\tag{11}
\end{equation}
sein, und $\xi$ ist eine Wurzel der Gleichung $f(x)=0$. Der Fundamentalsatz wird also dann bewiesen sein.

Beim Beweise von 5. machen wir Gebrauch von dem schon bewiesenen Satz (\S.\,32), dass eine reine Gleichung immer eine Wurzel hat.

Von den derivirten Functionen $f'(\alpha),f''(\alpha),\ldots$ können einige verschwinden; die letzte $f^{(n)}(\alpha)$, die gleich $\Pi(n)a_0$ ist, ist aber
%CONT>>
